% Generated draft: requires mathematical and source review.
\documentclass[11pt]{article}
\usepackage[margin=25mm]{geometry}
\usepackage{fontspec}
\setmainfont{Latin Modern Roman}
\usepackage{amsmath,amssymb,mathtools,unicode-math}
\setmathfont{Latin Modern Math}
\usepackage{enumitem,needspace,xurl,hyperref}
\setlength{\parindent}{0pt}
\setlength{\parskip}{5pt}
\setlength{\emergencystretch}{3em}
\newcommand{\N}{\mathbb{N}}
\newcommand{\Z}{\mathbb{Z}}
\newcommand{\Q}{\mathbb{Q}}
\newcommand{\R}{\mathbb{R}}
\newcommand{\C}{\mathbb{C}}
\newcommand{\F}{\mathbb{F}}
\newcommand{\A}{\mathbb{A}}
\newcommand{\PP}{\mathbb{P}}
\newcommand{\E}{\mathbb{E}}
\newcommand{\eps}{\varepsilon}
\newcommand{\dd}{\,\mathrm{d}}
\begin{document}
\section*{UnsolvedMath ID: 800605}
\textbf{Source problem code:} catalogue entry
\subsection{Definitions and mathematical statement}
The source asks whether \(\mathsf{Split}(\Lambda,\Lambda)\) is additive, but does not supply definitions of \(\Lambda\), additivity, or the class of spaces under consideration. The following makes a conventional interpretation explicit; these conventions require confirmation against the source.

Let \(T\) be a topological space and let \(\mathcal C\) be a specified class of subspaces of \(T\). For a subspace \(X\subseteq T\), define \(\Lambda(X)\) to consist of countably infinite families \(\mathcal U\) of open subsets of \(X\) such that \(X\notin\mathcal U\) and
\[
\forall x\in X,\qquad \bigl|\{U\in\mathcal U:x\in U\}\bigr|=\aleph_0.
\]
Such a family is called a large open cover. Covers here are families of distinct sets, rather than sequences with repetitions. Say that \(X\) satisfies \(\mathsf{Split}(\Lambda,\Lambda)\) if every \(\mathcal U\in\Lambda(X)\) has a partition
\[
\mathcal U=\mathcal U_0\mathbin{\dot\cup}\mathcal U_1,
\qquad \mathcal U_0,\mathcal U_1\in\Lambda(X).
\]
Under the interpretation that additivity means preservation under finite unions, the question is whether
\[
\forall X,Y\in\mathcal C,\quad
\bigl(X\text{ and }Y\text{ satisfy }\mathsf{Split}(\Lambda,\Lambda)\bigr)
\Longrightarrow
X\cup Y\text{ satisfies }\mathsf{Split}(\Lambda,\Lambda),
\]
where the union carries the subspace topology inherited from \(T\). Neither \(T\) nor \(\mathcal C\) is specified in the supplied excerpt. The source title concerns sets of reals, but the excerpt does not establish the exact domain. If additivity instead means preservation under countable unions, the intended assertion would replace \(X,Y\) by a sequence \((X_n)_{n\in\N}\) and the conclusion by the property for \(\bigcup_{n\in\N}X_n\). The supplied record does not resolve this distinction.

The excerpt also states: ``All problems below are settled for the properties which do not appear in them.'' This is retained as a source-specific qualification about the surrounding list; the omitted list and its property conventions are unavailable here.
\subsection{Short English statement}
Does the property that every large open cover can be partitioned into two large open covers persist under unions? The displayed formulation uses finite unions and countable covers, subject to confirmation of those conventions and the intended class of spaces. The registry labels the problem open; that label is supplied metadata, not an independently verified status.
\subsection{Sources}
\begin{itemize}
\item Local UnsolvedMath registry, permanent ID \(800605\), source code \texttt{catalogue entry}, AMR Open Problem Lists. Supplied upstream revision: \texttt{d371ee917ac9647fbb06cd6fc529deec6f3dda82}. The record attributes the problem to Tsaban, with proposed year \(2006\), and labels it open.
\item Tsaban, \emph{Selection Principles and Special Sets of Reals: Open Problems} (2006), Problem 33. Supplied source URL: \url{https://arxiv.org/abs/math/0609601}. The background records access on 2026-10-16 and extraction by \texttt{arxiv\_latex}; these are supplied provenance details, not a claim of independent source consultation.
\item The supplied attribution/context marker is \texttt{[split]}. Its bibliographic expansion is absent from the record.
\end{itemize}
\end{document}
